\chapter{The Original Carriers and Sixteen Charges}\label{ch:source}
\chaplead{The normalization of the color basis, Majoranas, fermion field and charge column fixes the factors used in every subsequent estimate.}

\section{Bosonic coordinates and color}
Use $X_{ai}=x_i^a$ for three color rows and nine spatial columns. Set
\[
 r^2=\tr(XX^{\mathsf T}),\qquad
 \Gram=XX^{\mathsf T},\qquad
 s=\tr\Gram,\quad u=\tfrac12((\tr\Gram)^2-\tr(\Gram^2)),\quad v=\det\Gram .
\]
The letter $u$ is an invariant in formulas involving $(s,u,v)$ and is a Hilbert vector in operator-domain formulas. The source uses zero-based indices for the finite types. We use one-based spatial and spin labels in displayed formulas.

Let $e_a=\sigma_a/\sqrt{2}$ be the normalized Hermitian Pauli color basis. Its trace pairing is orthonormal. For the source bracket $[E,F]_{\rm s}=-i(EF-FE)$, the structure coefficients are $f_{abc}=\sqrt{2}\,\epsilon_{abc}$ in the oriented Pauli order. Indeed $[\sigma_b,\sigma_c]=2i\sum_a\epsilon_{abc}\sigma_a$, and the two input factors $1/\sqrt{2}$ give $[e_b,e_c]_{\rm s}=\sum_a\epsilon_{abc}\sigma_a$. This convention fixes the crossover to the ordinary cross product.

\section{The fermionic fiber}
The real symmetric $16\times16$ matrices $\gamma_i$ obey
\[
 \gamma_i\gamma_j+\gamma_j\gamma_i=2\delta_{ij}I,\qquad
 \gamma^{ij}=\tfrac12(\gamma_i\gamma_j-\gamma_j\gamma_i).
\]
They act on a spatial spin label. The forty-eight original Hermitian fermions $\theta_{\alpha a}$ act on $\F$, with anticommutator
\[
 \{\theta_{\alpha a},\theta_{\beta b}\}=\delta_{\alpha\beta}\delta_{ab}I.
\]
It is convenient to use $t_{\alpha a}=\sqrt{2}\,\theta_{\alpha a}$, which square to $I$. In a paired mode with creation and annihilation $c_j^\dagger,c_j$ satisfying the complex CAR, the two normalized Majoranas are $c_j+c_j^\dagger$ and $i(c_j^\dagger-c_j)$ up to the oriented sign convention of the chosen pairs. The computational packing uses the actual source pair order and signs, which are part of its isometry receipt.

The Fock space is $\bigwedge^\bullet\C^{24}$ and $\dim_\C\F=2^{24}$. The multiplication algebra generated by all normalized Majoranas is $\operatorname{End}_\C\F$. One can see surjectivity directly: the number operators give the vacuum projection $\prod_j(1-c_j^\dagger c_j)$, and creation and annihilation words on its two sides give matrix units between occupation basis vectors. The dimension of the complex Clifford algebra for forty-eight generators is $2^{48}$, equal to the dimension of this endomorphism algebra. The full algebra has scalar commutant. Later uses on actual gauge states and accessible observables require their own constraints.

\evidence{\src{lean/OAI/MathematicalPhysics/BFSS/Core.lean}{Original source core}.
The theorem atlas records the carrier and original CAR results. The packed transport used for the trial is indexed in Chapter~\ref{ch:witness}.}

\section{Smooth core and gauge action}
Let $\Ltwo=L^2(\R^{27},dX;\F)$ with inner product conjugate-linear in its first argument. The compact smooth core $\Core=C_c^\infty(\R^{27};\F)$ embeds injectively in $\Ltwo$: a continuous function that vanishes almost everywhere with respect to full-support Lebesgue measure vanishes everywhere.

Write $R_g$ for the adjoint color matrix and $V_g$ for the actual fermionic unitary of the source gauge action. In these conventions the gauge action on sections is
\[
 (U_g f)(X)=V_g f(R_g^{-1}X).
\]
The adjoint color matrices lie in $\SO(3)$ and the adjoint map from $\SU(2)$ is onto. Orthogonality makes the bosonic change of variables measure preserving. The physical Hilbert space and core are
\[
 \E=\bigcap_g\ker(U_g-I),\qquad
 \PhysCore=\{f\in\Core:f(R_gX)=V_g f(X)\ \forall g,X\}.
\]
Boundedness of each $U_g$ makes the fixed space closed. For compact gauge group let $P_G f=\int U_gf\,dg$ with Haar probability measure. Averaging is contractive and fixes physical vectors. On compact smooth inputs the union of the rotated supports is bounded, and derivatives can be taken under the integral. Hence $P_G\Core\subset\PhysCore$. Applying $P_G$ to ambient smooth approximants proves density of the physical core in $\E$. The corresponding source Haar operation is used again at graph level.

\section{Nontriviality and the Gauss test}
An original nonzero gauge-invariant Fock vector multiplied by a nonzero smooth radial compact bump is a physical normalizable vector. The radial factor is invariant under the adjoint orthogonal color action. The norm is the product of its scalar $L^2$ norm and the fixed fiber norm, so positivity is immediate. The construction establishes that $\E$ is nonzero. Its differential energy is a separate calculation.

The source also has a historical maximal-domain Gauss characterization. On the ambient Hilbert space, the joint kernel of the actual Gauss generators at their maximal domains equals the fixed space of the genuine $\SU(2)$ action. That identification is stronger than the core assertion obtained by differentiating gauge invariance of smooth functions. Its historical receipt and exact domain interface are retained in the atlas.
\evidence{The historical radial construction and maximal Gauss equivalence are atlas G-02 and G-04. The original invariant fiber vacuum is G-03.}

\section{The literal differential operators}
For each $\alpha$ the source supercharge is
\begin{align}
 Q_\alpha f={}&-i\sum_{i,a,\beta}(\gamma_i)_{\alpha\beta}
           \theta_{\beta a}\,\partial_{ia} f \nonumber\\
 &+\frac12\sum_{i,j,a,b,c,\beta}f_{abc}(\gamma^{ij})_{\alpha\beta}
           \theta_{\beta a}x_i^b x_j^c f .
 \label{eq:sourcecharge}
\end{align}
The spatial pair in the second line is ordered. Equivalently, writing $X_i$ for the $i$th color column and using the normalized Majoranas gives the unordered-pair expression
\[
 Q_\alpha=-\frac{i}{\sqrt2}\sum_{i,a,\beta}
       (\gamma_i)_{\alpha\beta}t_{\beta a}\,\partial_{ia}
   +\sum_{i<j,a,\beta}(\gamma_i\gamma_j)_{\alpha\beta}
       (X_i\times X_j)_a\,t_{\beta a}.
\]
The $i<j$ pair occurs twice in the ordered sum, $f_{abc}$ supplies the factor $\sqrt2$, and $\theta=t/\sqrt2$. These factors explain the displayed coefficient of one for the polynomial term.

Each $Q_\alpha$ takes smooth compact inputs to smooth compact outputs. Its first-order coefficients have the form $-i$ times constant Hermitian matrices, and its zeroth-order field is pointwise Hermitian. Integration by parts gives $\ip{Q_\alpha f}{g}=\ip{f}{Q_\alpha g}$ on the ambient core. Gauge covariance makes the invariant core stable under the charges. For the derivative term this uses the orthogonal chain rule and the gauge action on the Majoranas. For the polynomial term it uses preservation of the color bracket under conjugation.

The source quadratic form is
\[
 q_{\rm core}[f]=\frac1{16}\sum_{\alpha=1}^{16}\norm{Q_\alpha f}^2 .
\]
This charge normalization carries through to the adjoint-square Hamiltonian. Formal symmetry and the finite sums are the starting data for the closed-operator work.
